\section{Discussion}\label{sec:discussion}

\subsection{Threats to Validity}\label{sec:threats}

\emph{Internal.} The same author identified the attackers, wrote the labelling rule and made all design choices. The labels were revised once after error analysis (the sixth DDoS source), which moved labels towards the model's predictions; we mitigated this by re-checking every flood victim over the whole file and by releasing the evidence. Timing comes from WSL2 on a laptop, where tcpreplay shares the CPU with the gateway. Run~1 and the one-core run used an earlier harness without the catch-up wait and per-window reset; all other runs used the final harness (Section~\ref{sec:testbed}). The validation simulation of blocking ignores the export delay (Section~\ref{sec:response}).

\emph{Construct.} Macro-F1 weighs tiny classes like large ones: Mirai has 10 and MITM-ARP 25 temporal test flows, and all Mirai test flows continue training connections. We therefore report per-class F1 with support and macro-F1 over classes with at least 50 test flows (0.814). False blocks are counted per host pair; scoring by the triggering flow gives identical counts (60 of 79 and 0 of 9).

\emph{External.} We use one dataset with one capture session per attack class, so the results measure generalisation within a session, not across networks or campaigns. Large captures were trimmed to their first $10^6$ packets, so a temporal split of such a file covers only the start of the capture session. The testbed is an emulated network on a laptop under a CPU quota, not a physical edge device; the gateway needs about 1\,GB of memory; and the live results hold only at the achieved replay rates (52--77\% of the original rate in the heaviest windows) and for floods that did not fill the capture buffer. A verdict needs the flow to end, and blocking acts on host pairs.

\emph{Conclusion.} The no-limit and half-core settings were run three times each, the one-core and blocking settings once. Macro-F1 and drops were stable across repeats, but flow counts and verdict delay varied with the achieved replay rate (Table~\ref{tab:repeats}). Flow-level bootstrap intervals are far too narrow (Section~\ref{sec:setup}). Under the host-pair bootstrap, the random- and temporal-split intervals overlap substantially, so the macro-F1 gap rests on a few host pairs and is not statistically robust. The SSH-to-RTSP change is a mechanism shown on three host pairs (Section~
ef{sec:splits}), a case study rather than a statistically robust effect. Statements such as ``no packets were lost'' refer only to the runs reported here.

\subsection{Lessons: Five Cheap Checks}\label{sec:lessons}

\begin{enumerate}
\item \textbf{Label from evidence, and price file-name labels.} Score any file-label model against evidence labels, and read confusions concentrated on one host pair as possible label errors. Background traffic was 91--98.5\% of several attack captures; file-name labels raised the true Benign FPR from 3.8\% to 71.8\%.
\item \textbf{Split by time within each capture, and read per-class F1 with support and host-pair intervals.} The random split hid a brute-force failure (recall 1.00 vs 0.05, three host pairs) caused by a change of target.
\item \textbf{Count boundary-spanning connections.} 36.1\% of temporal test flows continued training connections, including every Mirai test flow.
\item \textbf{Match live to offline flows, account for packets per window, and report achieved replay rates.} Equal scores can hide missing flows; packet accounting exposed a forwarding pitfall, and the per-flow match bounds how far live flows differ.
\item \textbf{Measure verdict delay against flow end under the target CPU budget, and score blocking in false blocks.} Report blocked attack traffic and the curve over $k$ per class, and choose the operating point on validation data. Delay exposed idle meters and an exporter defect that offline flushing hides; false-block counts exposed 60 benign pairs blocked on first verdicts and the loss of Mirai from $k=3$ on.
\end{enumerate}

\subsection{Conclusion}

We audited one flow-based IoT detector from its labels to live blocking. At each step a problem appeared that the previous step had hidden: file-name labels that turn background traffic into attacks, a random split that hides a change of brute-force target, live export problems that offline scores cannot show, and a per-flow false-positive rate that becomes dozens of wrongly blocked host pairs. Each was caught by a check that costs little compared with the experiment it protects. Future work includes repeated and longer live runs, early classification after the first packets of a flow to remove the dependence on the idle timeout, cross-dataset evaluation, and physical edge hardware.

\subsection{Artifact Availability}

Code, the evidence-based attacker definitions, configuration, result files and scripts that reproduce every number are available at \url{https://github.com/muhammadusamahoyrr/EdgeFlow-IDS} (MIT licence), including the labelling and extraction pipeline, models, tuning scripts, gateway, ARP watcher, Mininet/OVS topology and live harness. The CIC-IoT2023 pcaps are available from the dataset authors~\cite{neto2023ciciot2023}. The NFStream defect and its fix were reported as nfstream\#262~\cite{nfstream262}.
